\section{Conclusion et perspectives}

Corriger la récupération a apporté le gain le plus net et le mieux établi : sur le jeu de test, Hit@6 passe de
\NVoneTestHitsix~\% à \NVoneoneTestHitsix~\% et le temps par question baisse fortement. Sur l'agent, nos résultats sont plus
prudents. La planification systématique des requêtes ne bat pas nettement la question brute : c'est un résultat négatif,
obtenu en confrontant au test une variante qui brillait sur le développement. La combinaison d'un indice de code doux et
d'une boucle Évaluer → Reformuler atteint \NLoopLoopGHitsix~\% sur le développement contre \NLoopAtwoHitsix~\% pour A2, mais
ce gain repose sur peu de concepts et n'est pas confirmé. Le vérificateur de support détecte la plupart des affirmations
fausses construites par mutation, et un contrôle numérique sans LLM contre les PDF complète ce dispositif.

\paragraph{Perspectives.}
\begin{itemize}
  \item Rédiger le jeu de cent questions, écrire les attentes chiffrées, puis mesurer v1, v1.1 et v2 une seule fois.
  \item Brancher le vérificateur dans le graphe, avec une réparation au plus.
  \item Corriger les défauts de l'index relevés par l'audit (pied de page resté dans des \term{chunks}, sections).
  \item Étendre le contrôle de stabilité à l'arabe, à l'anglais et aux questions qui demandent une clarification.
  \item Faire valider un échantillon des étiquettes et des réponses par un juriste.
\end{itemize}
